\section{Methods}
\subsection{Source data and aligned cohort}
Structures and partial charges were taken from the pinned QMOF Figshare release, article 13147324, version 18, archive file 59573735.\cite{Rosen2021,Rosen2022,QMOFv18} No new electronic-structure calculations or experimental measurements were performed. The common finite-charge cohort contains 9,171 structures and 1,420,532 atoms, including 46,994 metal atoms in 9,165 structures. Atomic records were aligned by QMOF identifier and source atom index, with element identity and source site order checked. The common cohort is the intersection of available finite charge assignments, rather than a sample selected for a particular chemical effect. The structural and charge tables and graph configuration \nolinkurl{621045db53a892f1} were fixed before the conditional comparisons.

For an ordered pair of charge definitions $A$ and $B$, the target is
\begin{equation}
 \Delta_i^{AB}=q_i^A-q_i^B.
\end{equation}
We report DDEC6--CM5, DDEC6--Bader, and CM5--Bader in this order. These targets are algebraically dependent, since $\Delta^{\mathrm{DDEC6,Bader}}=\Delta^{\mathrm{DDEC6,CM5}}+\Delta^{\mathrm{CM5,Bader}}$; they are not three independent replications. Neither definition is treated as ground truth.

\subsection{Periodic local environments}
Structures were handled with pymatgen and ASE.\cite{Ong2013,Larsen2017} Metals follow \nolinkurl{pymatgen.core.Element.is_metal}. Nonmetal--nonmetal contacts were accepted between 0.40~\AA\ and 1.25 times the sum of ASE covalent radii. Metal-involving contacts used CrystalNN\cite{Zimmermann2020} without oxidation-state assignments and with electronegativity weighting disabled. A metal--nonmetal edge was accepted when selected by the metal center; a metal--metal edge was accepted when selected by either metal center. Every accepted edge was given an explicit reverse, retaining distinct periodic images of the same indexed neighbor. The complete parameter set and ambiguity definitions are given in the supporting information.

First-shell features describe the composition and distances of the selected neighborhood. Two-hop features count non-backtracking paths through that graph; they do not count unique atoms at graph distance two. Central-element information and the structural feature dictionary were fixed. No target charge, oxidation-state assignment, spin density, or electronic bond order entered graph selection. Diagnostic flags identify potentially ambiguous neighborhoods without declaring them chemically incorrect. Graph construction completed for every structure. The 72 zero-neighbor atoms and all flagged sites were retained.

\subsection{Conditional charge mappings and model comparison}
Within each training fold, an element-specific 12-bin piecewise charge mapping supplied a common prediction of $\Delta^{AB}$ from $q^A$. Training quantile bins, interpolation, endpoint extrapolation, and affine or global fallbacks are specified in the supporting information. Each competitor then predicted the residual to this common mapping. We compared an element residual mean, a flexible tree using central-element information and $q^A$, the same tree with first-shell features added, and the same tree with first-shell and two-hop features. Thus, only the charge-only to first-shell comparison isolates the first-shell contribution conditional on the source charge; the first-shell to two-hop comparison tests path information beyond that first shell. The charge-only to two-hop gain is the principal overall conditional comparison.

All tree fits used XGBoost\cite{Chen2016} with 100 estimators, maximum depth four, learning rate 0.08, minimum child weight 20, and $L_2$ penalty five. The algorithm, feature dictionary, training sample, and weights were matched across competitors. Training was capped at 160,000 atoms per fold through a deterministic structure--element sampling rule, with weights corrected for the retained sample. Every held-out atom was evaluated. No hyperparameter search, early stopping, or held-out model selection was performed. The source charge was permitted only in these explicitly conditional analyses; the target scheme's charge was never a predictor.

\subsection{Structure and identifier/composition validation}
Five original folds withheld entire structures. A second five-fold split linked records into connected components when they shared either a nonempty exact source MOFkey\cite{Bucior2019} or an exact source reduced formula. This created 8,347 components, including 1,414 records in nonsingleton components. Components were assigned intact to folds containing 1,835, 1,834, 1,834, 1,834, and 1,834 structures. All calibrations and residual models were newly fitted for this split. MOFkey-present and nonsingleton records were also evaluated separately. This tests the specified identifier/composition dependence; it does not establish separation of every chemical family or unseen topology.

The principal error is the structure-balanced mean absolute error,
\begin{equation}
 \mathrm{MAE}_{\mathrm{str}}=\frac{1}{N}\sum_{s=1}^{N}
 \frac{1}{n_s}\sum_{i\in s}\left|\Delta_i^{AB}-\widehat{\Delta}_i^{AB}\right|,
\end{equation}
where $n_s$ is the number of evaluated atoms in structure $s$ for the stated population. Atom-weighted errors are reported separately. Positive gain denotes baseline MAE minus structural-model MAE on the same held-out atoms. Intervals use 1,000 paired bootstrap draws over structures for the original folds and over components for the blocked folds. Component resampling retains sums of per-structure losses and structure counts, preserving the structure-balanced estimand. These pointwise intervals condition on the fitted models and do not include uncertainty from retraining or unrecognized broader families.

\subsection{Chemical atlas and graph sensitivity}
For the chemical atlas, the flexible charge-only held-out residual is $r_i=\Delta_i^{AB}-\widehat{\Delta}_{i,\mathrm{charge}}^{AB}$. We separately summarize $\Delta_i^{AB}$, $|\Delta_i^{AB}|$, $r_i$, $|r_i|$, and the within-structure/environment standard deviation of $r_i$. Same-element environment contrasts are paired within structures before averaging. Six fixed contrasts compare O-attached or N-attached carbon with strictly C/H-only carbon, O-bound or N-bound monocoordinate H with C-bound H, and metal-adjacent O or N with their non-metal-adjacent counterparts. N-attached carbon excludes O attachment. Compact atlas entries require at least 30 structures and 30 represented validation units; the full tables retain lower-support classes.

Graph sensitivity used a fixed 96-structure subset containing 16,165 atoms. The nonmetal radius scale was changed to 1.20 and 1.30 while accepted metal-involving edges were held fixed. Descriptors, environment membership, and paired contrasts were recomputed with the same held-out residuals, comparing common paired structures. Separate, earlier fixed-model perturbations of metal contacts are reported with their original model population and conditioning; they are not refits of the conditional models. The graph is an operational environment definition, not a universal chemical bond assignment.
